\documentclass{article}
% Source: Topic_08_Teacher_Edition.pdf, PDF pages 27-39 (printed pages 404A-412B)
\usepackage{amsmath}
\usepackage{amssymb}
\usepackage{graphicx}
\usepackage{tikz}
\usepackage{pgfplots}
\pgfplotsset{compat=1.18}
\usepackage{booktabs}
\usepackage{xcolor}

\newenvironment{lesson-meta}{}{}        % lesson code, title, objective, EQ, vocab
\newenvironment{item-analysis}{}{}      % Savvas's Example × DOK table
\newenvironment{model-discuss}{}{}      % Launch opener
\newenvironment{concept-box}{}{}        % in-lesson concept reference
\newenvironment{concept-summary}{}{}    % end-of-lesson summary
\newenvironment{example}[1]{}{}         % #1 = example number
\newenvironment{tryit}[1]{}{}           % #1 = tryit number
\newenvironment{practice}[2]{}{}        % #1 = practice number, #2 = DOK
\newenvironment{te-addendum}[2]{}{}     % #1 = anchor example, #2 = type
\newcommand{\answer}[1]{}               % answer key, inside any env
\newcommand{\placeholder}[2]{}          % #1 = type, #2 = description
\newcommand{\uncertain}[2]{#1}          % #1 = best guess, #2 = alternatives
\newcommand{\te}[1]{}                   % teacher-only inline annotation

\begin{document}
% Source page 27: 404A Lesson Overview
\begin{lesson-meta}
lesson: 8-2
title: Law of Sines and Law of Cosines
standards: none printed
practices: none printed
objective: I can use the Law of Sines and the Law of Cosines to solve for unknown angles and sides of a triangle.

essential question: How can you use the sine and cosine functions with non-right triangles?

\textbf{VOCABULARY}
\item Law of Cosines
\item Law of Sines
\textbf{TOPIC VOCABULARY}
\te{No topic vocabulary list is printed on these pages.}
\begin{tabular}{ll}
Lesson Code & 8-2 \\
Title & Law of Sines and Law of Cosines \\
Standards & none printed \\
I CAN Objective & I can use the Law of Sines and the Law of Cosines to solve for unknown angles and sides of a triangle. \\
Essential Question & How can you use the sine and cosine functions with non-right triangles? \\
Lesson Vocabulary & Law of Cosines; Law of Sines \\
Topic Vocabulary & none printed \\
\end{tabular}
\end{lesson-meta}
\begin{te-addendum}{0}{GeneralizeETP}
\textbf{Lesson Overview: FOCUS}
Objective. Students will be able to:
\item Derive the Law of Sines and the Law of Cosines.
\item Use the Law of Cosines and the Law of Sines to find unknown angles and sides of non-right triangles.
Essential Understanding. The Law of Sines or Law of Cosines can be used to find unknown measurements of a non-right triangle.
\textbf{COHERENCE}
Previously in this topic students:
\item Used sine and cosine functions to solve problems and to find unknown measurements of right triangles.
In this lesson, students:
\item Use the Law of Sines and the Law of Cosines to solve problems and to find unknown measurements of non-right triangles.
Increasing Coherence This lesson is not necessarily expected to be addressed on high stakes assessments.
Later in this topic, students will:
\item Use trigonometric identities, including sines and cosines, to verify relationships between trigonometric functions.
\textbf{RIGOR}
This lesson emphasizes a blend of conceptual understanding and application.
\item Students understand that the Law of Sines and the Law of Cosines extend the use of trigonometric functions to finding unknown measurements of non-right triangles.
\item Students apply the Law of Sines and the Law of Cosines to solve real-world problems, such as finding the angle between two points in a triangular path.
\textbf{Mathematics Overview}
Students prove the Law of Cosines and the Law of Sines. Then they use these Laws to solve mathematical and real-world problems that involve finding angle measures and side lengths of triangles that are not right triangles.
\textbf{Applying Math Practices}
Reason Abstractly and Quantitatively
Students exhibit flexibility in their use of the Law of Sines and Law of Cosines when they solve a problem using one angle and then consider whether they could have solved the problem using a different angle instead.
Look For and Make Use of Structure
Students apply general rules to a specific situation when they recognize that in order to apply the Law of Sines, they need to know the measure of one angle and the length of the side opposite that angle.
\te{Program Resources. SavvasRealize.com.}
\end{te-addendum}
\begin{te-addendum}{0}{ELLAddendum}
\textbf{Vocabulary Builder}
Glossary is available at SavvasRealize.com.
\begin{tabular}{lll}
 & English & Spanish \\
REVIEW VOCABULARY & Pythagorean Theorem & Teorema de Pitágoras \\
NEW VOCABULARY & Law of Cosines & Ley de cosenos \\
 & Law of Sines & Ley de senos \\
\end{tabular}
VOCABULARY ACTIVITY
Review the Pythagorean Theorem using a right triangle. Then describe the Law of Sines and the Law of Cosines using a non-right triangle. Have students match each of the equations to the law or theorem it describes.
1. (\frac{\sin A}{a}=\frac{\sin B}{b}=\frac{\sin C}{c})
2. (a^2+b^2=c^2)
3. (a^2=b^2+c^2-2bc(\cos A))
a. Law of Cosines
b. Law of Sines
c. Pythagorean Theorem
\answer{1--b; 2--c; 3--a.}
\end{te-addendum}
\begin{te-addendum}{0}{LearnTogether}
\textbf{Student Companion}
Students can do their work for the lesson in their Student Companion or on Savvas Realize.
% FIG 27 932 673 1087 853
\placeholder{illustration}{Student Companion cover with glowing colored loops on a dark background; enVision Algebra 2, SAVVAS.}
\end{te-addendum}
% Source page 28: 404B Explore
\begin{te-addendum}{0}{LearnTogether}
\textbf{MODEL \& DISCUSS}
INSTRUCTIONAL FOCUS Students revisit their understanding of using trigonometric functions to find missing side lengths of a right triangle. This leads students to understand and prove the Law of Sines and Law of Cosines. Students use these laws to solve for unknown angles and sides of non-right triangles.
STUDENT COMPANION Students can complete the Model \& Discuss activity in their Student Companion.
\te{Activity. Critique \& Explain is available at SavvasRealize.com. Printed Critique \& Explain; likely Model \& Discuss.}
\end{te-addendum}
\begin{te-addendum}{0}{GeneralizeETP}
\textbf{Before: WHOLE CLASS}
Implement Tasks that Promote Reasoning and Problem Solving ETP
Q: How can you create a drawing of the mound using right triangles?
\answer{Represent a vertical cross-section of the cone through its vertex. Then the slant height and diameter represent 2 sides of the triangle, and the included angle is the angle from the ground to the top of the mound. Then create two right triangles by drawing the height of the triangle from the vertex to the side representing the ground.}
\textbf{During: SMALL GROUP}
Support Productive Struggle in Learning Mathematics ETP
Q: What are two ways that you can estimate the height of the mound?
\answer{Using the given angle, apply a trigonometric function; Measure the height and one of the labeled sides with a ruler. Use proportional reasoning to estimate the height based on your measurement.}
For Early Finishers
Q: Draw a rectangle with a width of 24 in. and a diagonal that creates a 52° angle with the width. Use the given angle and trigonometric functions to find the length of both the diagonal and the rectangle.
\answer{Check students' diagrams. By drawing the diagonal, you create two right triangles, both of which have a 52° angle formed by the hypotenuse and the width of the rectangle. The diagonal of the rectangle has a length of approximately 39 in. and the length of the rectangle is approximately 31 in.}
\textbf{After: WHOLE CLASS}
Facilitate Meaningful Mathematical Discourse ETP
Q: How can you apply trigonometric functions to rectangles or triangles that are not right triangles?
\answer{In a rectangle, draw a diagonal and create two right triangles. In a triangle that is not a right triangle, draw an altitude perpendicular to the base and create two right triangles.}
\end{te-addendum}
\begin{model-discuss}
\textbf{MODEL \& DISCUSS}
A biologist measures the slant height of a conical termite mound to be about 32 ft. The angle from the ground to the top of the mound is 51°. The base of the mound has a diameter of about 40 ft.
A. Draw a model to help the biologist.
\te{Digital version: Use the pen tool to draw a model to help the biologist. Go back; ESSENTIAL QUESTION; Enter your answer; SOLUTION.}
\answer{A. See sample student work.}
% FIG 28 713 822 920 996
\placeholder{diagram}{Green filled triangle with gray outline, horizontal base labeled 40 ft., apex above and left of center, right slanted side 32 ft., arc at lower right labeled 51°.}
B. Make Sense and Persevere What is the height of the mound?
\answer{B. Draw an altitude, h, from the top of the triangle to the base. This equation can be used to find the height of the mound: (\sin51=\frac{h}{32}). The height is about 25 ft.}
\end{model-discuss}
\begin{te-addendum}{0}{HabitsOfMind}
Use with MODEL \& DISCUSS
Reason How did you decide which trigonometric function to use to solve the problem?
\answer{The information given includes the hypotenuse and adjacent angle, and the question asks to find the length of the side opposite the angle. This suggests using sine.}
\end{te-addendum}
% Source page 29: 404 Understand & Apply
\begin{te-addendum}{0}{LearnTogether}
\textbf{INTRODUCE THE ESSENTIAL QUESTION}
Establish Mathematics Goals to Focus Learning ETP
Students learn the Law of Sines and the Law of Cosines and how they are derived. Students also learn to apply the Law of Sines and Law of Cosines to solve for unknown angles and sides of a non-right triangle.
\te{Examples are available at SavvasRealize.com.}
\end{te-addendum}
\begin{concept-box}
\textbf{Law of Sines and Law of Cosines}
The Law of Sines and the Law of Cosines allow you to apply trigonometric functions to non-right triangles. Given (\triangle ABC), with angles A, B, and C and opposite-side lengths a, b, and c:
Law of Sines: (\frac{\sin A}{a}=\frac{\sin B}{b}=\frac{\sin C}{c})
Law of Cosines:
(a^2=b^2+c^2-2bc(\cos A))
(b^2=a^2+c^2-2ac(\cos B))
(c^2=a^2+b^2-2ab(\cos C))
% FIG 29 966 292 1120 357
\placeholder{diagram}{Black triangle ABC, horizontal base A left B right, C above and left of center. Red side labels b on AC, a on CB, c on AB.}
\end{concept-box}
\begin{te-addendum}{1}{PurposefulQuestions}
\textbf{Prove the Law of Sines}
Pose Purposeful Questions ETP
Q: What do you notice about (\triangle ABC) that is different from the triangles used in trigonometric functions?
\answer{(\triangle ABC) is not a right triangle.}
Q: How is drawing an altitude helpful?
\answer{Drawing an altitude creates two right triangles so you can apply the trigonometric functions.}
Q: How does isolating x in both sine functions help prove the Law of Sines?
\answer{Isolating x shows that x is equal to both b sin A and a sin B; therefore b sin A = a sin B.}
\end{te-addendum}
\begin{example}{1}
\textbf{Prove the Law of Sines}
How can you derive the Law of Sines?
Step 1 Draw (\triangle ABC) with an altitude from C to side c with length x.
% FIG 29 986 439 1110 495
\placeholder{diagram}{Black triangle ABC, base A left B right, C above. Red side labels b on AC, a on CB, c on AB; red vertical altitude x from C to AB with right-angle square at foot.}
\te{STUDY TIP Drawing an altitude allows you to create two right triangles and apply trigonometric functions.}
Step 2 Write the sine function for angles A and B.
(\sin A=\frac xb\text{ or }x=b\sin A)
(\sin B=\frac xa\text{ or }x=a\sin B)
Step 3 Since you have two statements that are both equal to x, you can set them equal to each other.
(b\sin A=a\sin B)
Step 4 Divide both sides of the equation by ab.
(\frac{\sin A}{a}=\frac{\sin B}{b})
The ratio of the sine of an angle to its opposite side is the same for all angles in the same triangle.
\answer{(\frac{\sin A}{a}=\frac{\sin B}{b})}
\end{example}
\begin{tryit}{1}
How can you derive the Law of Sines for angles B and C?
\answer{(\sin C'=\sin C=\frac yb); (\sin B=\frac yc); (y=b\sin C); (y=c\sin B); (c\sin B=b\sin C); (\frac{\sin B}{b}=\frac{\sin C}{c}).}
% FIG 29 312 651 604 805
\placeholder{diagram}{Black obtuse triangle ABC, horizontal base AB labeled red c, A left B right, C above and left of center. AC red b, CB red a. Dashed extension of BC past C up left to X; dashed altitude AX red y, right-angle square at X. Exterior angle at C marked C'.}
\end{tryit}
\begin{te-addendum}{5}{PurposefulQuestions}
\textbf{Additional Example 5}
A bicycle ramp is 10 ft long, its base is 13 ft, and the end has a slant height of 5 ft. What is the angle from the base to the top of the ramp?
% FIG 29 389 1006 510 1050
\placeholder{diagram}{Black scalene triangle with horizontal base 13, upper vertex near right, left sloping side 10, right side 5; angle A at left base in red.}
Since you know the lengths of the three sides of the triangle, you can use the Law of Cosines to solve for the angle of the ramp.
(a^2=b^2+c^2-2bc(\cos A))
(5^2=10^2+13^2-2(10)(13)(\cos A))
(-244=-260(\cos A))
(\frac{244}{260}=\cos A)
(\cos^{-1}(\frac{244}{260})=A)
(\cos^{-1}(0.938)\approx A)
(20.21°\approx A)
\answer{The angle from the base to the ramp is approximately 20.21°.}
Example 5 Help students transition to using the Law of Cosines to find an angle measurement.
Q: After you isolate the cosine function, what final step must be applied to find the angle measurement?
\answer{The inverse cosine function must be applied.}
\te{Go back; ESSENTIAL QUESTION; SOLUTION. Additional Examples are available at SavvasRealize.com.}
\end{te-addendum}
\begin{te-addendum}{6}{PurposefulQuestions}
\textbf{Additional Example 6}
The flight of a drone forms a triangle. If you fly the drone from home to point A and then fly to point B, how far from home is the drone at point B? What is the measure of the angle formed between the flight from point A to point B and the distance from point B to home?
% FIG 29 932 1002 1059 1110
\placeholder{diagram}{Triangle home-A-B, horizontal base B left A right 11 mi, top home marked blue dot, A and B black dots. Home-A black side 9 mi, home-B blue side. Arc at A labeled 51°.}
(a^2=b^2+c^2-2bc(\cos A))
(a^2=9^2+11^2-2(9)(11)(\cos51°))
\answer{(a\approx8.8), so the drone is about 8.8 miles from home.}
(\frac{\sin A}{a}=\frac{\sin B}{b})
(\frac{\sin51°}{8.8}=\frac{\sin B}{9})
\answer{The angle formed is approximately 52.6°.}
Example 6 Students use the Law of Cosines and the Law of Sines to answer a more complex multi-step problem.
Q: When setting up a multi-step problem that requires the use of both laws, what is the initial step?
\answer{The initial step is to designate the angle and side upon which the Law of Cosines will be based.}
\te{Go back; ESSENTIAL QUESTION; SOLUTION.}
\end{te-addendum}
% Source page 30: 405 Use the Law of Sines
\begin{te-addendum}{2}{PurposefulQuestions}
\textbf{Use the Law of Sines}
Pose Purposeful Questions ETP
PART A
Q: Under what circumstances can you use the Law of Sines?
\answer{In order to use the Law of Sines, you need to know either the measure of 2 angles and the length on a side opposite a known angle, or two side lengths and the measure of one angle opposite a known side.}
Q: When using the Law of Sines, how do you isolate the sine function to solve for an angle measurement?
\answer{Multiply each side of the equation by its denominator.}
PART B
Q: Which angle do you use when writing the equation to find the distance between the two people?
\answer{(\angle K); The distance is opposite (\angle K); therefore you need to calculate (m\angle K) by subtracting ((m\angle J+m\angle L)) from 180°, and then use the result in the equation.}
Q: Do you use the inverse sine function to find the distance between the two people?
\answer{No, the length is found by evaluating the ratio; inverse functions are used when solving for angle measurements.}
\te{Examples are available at SavvasRealize.com.}
\end{te-addendum}
\begin{te-addendum}{2}{ELLAddendum}
\textbf{English Language Learners} (Use with EXAMPLE 2)
SPEAKING BEGINNING Discuss with students how one definition for obstacle is an object that makes it difficult for you to go where you want to go because it is in your way.
Q: Describe some obstacles you might encounter on an obstacle course.
\answer{a mud pit, a rock wall, a tunnel you have to crawl through}
Q: Are obstacles always physical objects?
\answer{No, uncorrected vision could be an obstacle to reading. Being short could be an obstacle to driving.}
READING DEVELOPING Have students read Part A of the example along with you. Then have them read the definitions for trail from a dictionary: (verb) to drag, follow, or hang down; (noun) path, track, or route. Note that path and trail are synonyms.
Q: What part of speech is trail as it is used in the example?
\answer{noun}
Q: The word path is also used throughout the example. Does path have the same meaning as trail each time?
\answer{Yes; trail and path mean route or track.}
WRITING EXPANDING Research zeppelins online. Have students write the definition and a paragraph about zeppelins in their journals. Then have them fill in the blanks to summarize zeppelins.
Q: A zeppelin is a large aircraft without blank that blank because it is filled with gas and has a rigid frame inside its body to help it keep its blank.
\answer{wings, floats, shape}
Q: Zeppelins were used as blank during World War I. Their use was discontinued after the famous blank disaster in 1937.
\answer{bombers, Hindenburg}
\end{te-addendum}
\begin{example}{2}
\textbf{Use the Law of Sines}
\te{APPLICATION}
A. Naeem is walking along an obstacle course that begins at point X. He starts going straight on the path. It then takes a wide right turn onto a second path (at point Y) and continues for 1,000 feet before turning on to the final path (at point Z). What is the angle between the second path and the final path (point Z)? Round to the nearest tenth of a degree.
% FIG 30 985 237 1181 398
\placeholder{illustration}{Overhead obstacle course among trees, triangular paths X lower left, Y upper left, Z upper right. XY 1,200 ft, YZ 1,000 ft; X angle 32°, Z marked ?. Black arrows along XY then YZ then ZX; rope bridges and person on ZX. Not drawn to scale.}
(\frac{\sin A}{a}=\frac{\sin B}{b})
(\frac{\sin32°}{1,000}=\frac{\sin Z}{1,200})
(\frac{6\sin32°}{5}=\sin Z)
(0.636\approx\sin Z)
(\sin^{-1}(0.636)\approx\angle Z)
(\angle Z\approx39.5°)
\te{Callout: Substitute the values into the equation and begin to solve for Z. USE STRUCTURE In order to apply the Law of Sines, you must know the measure of one angle and the length of the side opposite the angle.}
One angle with a sine of 0.636 is (\sin^{-1}(0.636)\approx39.5°). The angle (180°-39.5°=140.5°) also has a sine of 0.636. Since Z is acute, (m\angle Z\approx39.5°).
\answer{The angle between the second path and the final path is about 39.5°.}
B. Amaya is flying her zeppelin balloon. The string is 90 ft long, and the angle of elevation to the balloon from the ground is 60°. Across the park, Rochelle is watching the balloon, which is between Rochelle and Amaya. The angle of elevation from Rochelle's feet to the balloon is 75°. How far apart are Amaya and Rochelle standing? Round to the nearest tenth of a foot.
% FIG 30 993 562 1183 685
\placeholder{illustration}{Park with pond, trees and balloon above. Black triangle JKL connects Rochelle J lower left, balloon K above, Amaya L lower right; JL horizontal, KL 90 ft. Angle J 75°, L 60°.}
To use the Law of Sines to find JL, you need to know the measure of the angle opposite JL.
Find (m\angle K) using the measures of the other two angles in the triangle.
(m\angle K=180°-75°-60°=45°)
\te{REASON Could you have used (\angle L) to solve the problem instead? CONTINUED ON THE NEXT PAGE.}
% Source page 31: 406 Example 2 continued
(\frac{\sin A}{a}=\frac{\sin B}{b})
(\frac{\sin45°}{k}=\frac{\sin75°}{90})
(\frac{90\sin45°}{\sin75°}=k)
(k\approx65.9\text{ ft})
\te{Callout: Solve for k by cross-multiplication and then divide by sin 75°.}
\answer{Amaya and Rochelle are standing about 65.9 ft apart.}
\end{example}
\begin{tryit}{2}
In (\triangle NPQ), (m\angle N=105°), (n=12), and (p=10).
a. To the nearest degree, what is (m\angle Q)?
b. What is the length of side q? Round to the nearest tenth of a unit.
\answer{a. 21.4° b. 4.5 units}
\te{Printed part a asks nearest degree but key gives 21.4°; likely 21° at requested precision.}
\end{tryit}
\begin{te-addendum}{2}{ElicitEvidence}
Elicit and Use Evidence of Student Thinking ETP
Q: When using the Law of Sines, what information do you need to find before you can solve for (m\angle Q) and the length of q?
\answer{(m\angle P)}
\end{te-addendum}
\begin{te-addendum}{2}{CommonError}
Try It! 2 Some students may incorrectly use the information given, (\frac{\sin105}{12}=\frac{\sin Q}{10}). Encourage students to write the Law of Sines for the triangle they are currently working with in each problem to stay organized: (\frac{\sin N}{n}=\frac{\sin P}{p}=\frac{\sin Q}{q}).
\end{te-addendum}
\begin{te-addendum}{3}{GeneralizeETP}
\textbf{Understand the Ambiguous Case}
Use and Connect Mathematical Representations ETP
Q: What does A refer to in sin A?
\answer{A refers to the measure of angle A. Because this value is always an angle measure, convention says you only write the angle name.}
Q: How is it possible for (m\angle E) to have two possible values if the angle looks small?
\answer{When the triangle is defined, nothing specifies that (\angle E) is acute. Drawings are not always to scale, so you must work out the measurement mathematically.}
Q: What causes the ambiguous case?
\answer{The sine function is positive for (0°<x<180°). If one acute angle is given, then two triangles may be possible.}
Q: After finding one possible angle, how do you find the other option?
\answer{The two angle measure possibilities are supplementary, so subtract the known angle from 180°.}
\end{te-addendum}
\begin{te-addendum}{3}{RtIExtend}
\textbf{Extend Student Thinking}
USE WITH EXAMPLE 3 Students explore the ambiguous case by finding all unknown angle measures and side lengths in both possible triangles.
\item In (\triangle ABC), (c=11), (b=8), and the (m\angle B=23°). Find the possible measurements of (\angle C), (\angle A), and a.
Q: Explain how to find the first set of measurements for (\angle C), (\angle A), and a.
\answer{Use the Law of Sines, (\frac{\sin23°}{8}=\frac{\sin C}{11}), to find (m\angle C\approx32.5°). Then use the fact that the sum of the angles of a triangle is 180° to find (m\angle A): (m\angle A=180-23-32.5=124.5°). Use the Law of Sines a second time to find BC: (\frac{\sin23°}{8}=\frac{\sin124.5°}{a}), so (a\approx16.87).}
Q: Explain how to find the second set of measurements for (\angle C), (\angle A), and a.
\answer{Use the unit circle to see that a 147.5°-angle has the same sine as a 32.5°-angle. So, (m\angle C) could be 147.5°. Then use the fact that the sum of the angles of a triangle is 180° to find (m\angle A): (m\angle A=180-23-147.5=9.5°). Use the Law of Sines to find a: (\frac{\sin23°}{8}=\frac{\sin9.5°}{a}), so (a\approx3.38).}
\end{te-addendum}
\begin{example}{3}
\textbf{Understand the Ambiguous Case}
\te{CONCEPTUAL UNDERSTANDING}
In (\triangle CDE), (CD=8), (DE=6), and (m\angle C=30°). What is (m\angle E) in (\triangle CDE)? Round to the nearest degree.
Start by drawing a sketch of the given information. It is not necessary to try to make your sketch to scale as you will not be relying on measurements of the sketch to find the missing information.
% FIG 31 974 469 1114 533
\placeholder{diagram}{Black triangle with E lower left, C lower right, D above and left of center. Red lengths ED 6, DC 8; angle C 30°.}
Because you know one side and the angle opposite that side, you can use the Law of Sines to find a second angle.
(\frac{\sin A}{a}=\frac{\sin B}{b})
(\frac{\sin30°}{6}=\frac{\sin E}{8})
(\frac{8\sin30°}{6}=\sin E)
(\sin^{-1}(\frac{4\sin30°}{3})=\angle E)
(\angle E\approx42°)
\te{Callout: Substitute and then isolate the variable.}
So you could conclude that the angles in the triangle measure 30°, 42°, and 108°. However, there is another angle which has the same sine value as the 42° angle.
Using the unit circle, you can see that an angle with measure 138° has the same sine value.
% FIG 31 996 684 1115 806
\placeholder{graph}{Red unit circle on black double-arrowed x,y axes, O origin. Square grid spacing .25, window -1.5 to 1.5 on both axes; x labels -1,-0.5,0.5,1 and y labels -1,-0.5,1. Blue radii to filled blue quadrant I and II points at angles 42° and 138°. Green vertical segments to x axis labeled sin 42° and sin 138°. Red 42° arc and green 138° arc from positive x axis.}
\te{LOOK FOR RELATIONSHIPS If a triangle has one obtuse angle given, there is only one possible triangle. If one acute angle is given, then two triangles may be possible. CONTINUED ON THE NEXT PAGE.}
% Source page 32: 407 Example 3 continued
So the angles in the triangle could also measure 30°, 138°, and 12°. The triangle would look quite different from the original sketch, but it does match all of the information given in the problem.
% FIG 32 1005 221 1144 289
\placeholder{diagram}{Triangle EDC with horizontal base E left C right, D above left; ED 6, DC 8 solid, angle C 30°. Dashed DE' from D to E' on base near C; obtuse angle at E' 138°. Original E labeled E; new base point labeled E.}
There are two possible triangles that could have the given information, one where (m\angle E=42°) and one where (m\angle E=138°).
\te{STUDY TIP Regardless of what information a problem is asking for or how it is being solved, do not rely on the appearance of an angle to guide your answers. Always work measurements out mathematically.}
\answer{(m\angle E=42°) or (m\angle E=138°).}
\end{example}
\begin{tryit}{3}
In (\triangle ABC), (m\angle A=30°), (a=5), and (b=8). Find (m\angle B). How many possible triangles are there?
\answer{(m\angle B=53°) or 127°. There are 2 possible triangles.}
\end{tryit}
\begin{te-addendum}{3}{ElicitEvidence}
Elicit and Use Evidence of Student Thinking ETP
Q: Do the given side lengths of the triangle change since there are two different triangles possible?
\answer{No; the given lengths remain the same. The side length that is not given is different in each triangle because its opposite angle is different for each triangle.}
\end{te-addendum}
\begin{te-addendum}{3}{HabitsOfMind}
Use with EXAMPLES 1, 2, \& 3
Use Structure In (\triangle ABC), (m\angle A=50°), (b=6), and (c=3). Gregory says that the Law of Sines does not allow you to find side length a. Is he correct? If so, what would he need to know in order to be able to find it?
\answer{Yes, he would need to know either of the other angle measures in order to use the Law of Sines.}
\end{te-addendum}
\begin{te-addendum}{4}{PurposefulQuestions}
\textbf{Prove the Law of Cosines}
Pose Purposeful Questions ETP
Q: What do you notice about the Law of Cosines?
\answer{It is based on the Pythagorean Theorem and uses right triangles.}
Q: How do you start proving the Law of Cosines?
\answer{Begin by drawing an altitude to create two right triangles, where the hypotenuse of one right triangle is opposite the angle you are using to prove the Law of Cosines.}
\end{te-addendum}
\begin{example}{4}
\textbf{Prove the Law of Cosines}
Prove the Law of Cosines (a^2=b^2+c^2-2bc(\cos A)) for an acute angle A.
% FIG 32 839 430 961 493
\placeholder{diagram}{Black triangle ABC with horizontal base A left B right, C above; AC b, CB a; dimension arrow under AB c. Orange altitude CD with right-angle square at D on AB, green x label. AD red y, DB blue z.}
Step 1 The Law of Cosines resembles the Pythagorean Theorem. To prove the Law of Cosines, start by thinking of a right triangle in which a is the hypotenuse.
The drawing shows (\triangle ABC) with an altitude of length x that divides length c into y and z at point D, so that two right triangles are formed.
Step 2 Since z is in the right triangle with hypotenuse a, find an expression for z using b, c, and (\angle A), the variables on the right side of the Law of Cosines equation.
(\cos A=\frac yb\text{ or }y=b\cos A)
(c=y+z\text{ or }z=c-y)
(z=c-b\cos A)
\te{Callout: Substitute b cos A for y.}
Step 3 Now you need an expression for x in terms of b, c, and (\angle A). You can get one by finding sin A.
(\sin A=\frac xb\text{ or }x=b\sin A)
\te{CONTINUED ON THE NEXT PAGE}
% Source page 33: 408 Example 4 continued
Step 4 Use the Pythagorean Theorem for (\triangle BCD), and substitute to obtain an equation relating a, b, c, and (\angle A). Simplify.
(a^2=x^2+z^2)
(a^2=(b\sin A)^2+(c-y)^2)
(a^2=(b\sin A)^2+c^2-2cy+y^2)
(a^2=b^2\sin^2 A+c^2-2bc(\cos A)+b^2\cos^2 A)
(a^2=b^2(\sin^2 A+\cos^2 A)+c^2-2bc(\cos A))
(a^2=b^2+c^2-2bc(\cos A))
\te{Callouts: Substitute (y=b\cos A). Rearrange and factor (b^2). Remember, (\sin^2x+\cos^2x=1). COMMON ERROR Be careful to square both quantities correctly. The first quantity is a monomial, while the second is a binomial.}
This result is the Law of Cosines for (\angle A).
\answer{(a^2=b^2+c^2-2bc(\cos A))}
\end{example}
\begin{tryit}{4}
a. How can you derive the Law of Cosines for an obtuse angle C?
% FIG 32 180 849 404 1013
\placeholder{diagram}{Black obtuse triangle BCA, B lower left, C on horizontal base near center, A upper right. BC red a, CA red b, BA red c. Dashed extension CD right, red x; dashed vertical altitude AD red h; right-angle square at D. Exterior angle ACD marked C'.}
\answer{a. (\triangle ABD:c^2=(a+x)^2+h^2=a^2+2ax+x^2+h^2). (\triangle ACD:b^2=x^2+h^2), so (c^2=a^2+2ax+b^2); (\cos C=-\cos C'=-\frac xb), so (x=-b\cos C). Therefore, (c^2=a^2-2ab\cos C+b^2=a^2+b^2-2ab\cos C).}
b. How does the equation compare to the equation for an acute angle?
\answer{b. The equation is the same for both acute and obtuse angles.}
\end{tryit}
\begin{te-addendum}{5}{PurposefulQuestions}
\textbf{Use the Law of Cosines}
Pose Purposeful Questions ETP
Q: When you determine which side of the triangle is represented by a, does it matter which sides are designated as b and c?
\answer{No; Addition and multiplication are commutative. The angle designated as A and its opposite side must be identified correctly.}
Q: When using the Law of Cosines, why is using technology useful?
\answer{Since the Law of Cosines is based on the Pythagorean Theorem, the square root of non-perfect squares need to be approximated.}
\te{Examples are available at SavvasRealize.com.}
\end{te-addendum}
\begin{example}{5}
\textbf{Use the Law of Cosines}
What is t? Round to the nearest tenth.
% FIG 33 797 514 913 587
\placeholder{diagram}{Black triangle TUV, horizontal base T left V right 17 in red, U above left, TU 9, UV t, angle T 65° in red.}
\begin{tabular}{ll}
(a^2=b^2+c^2-2bc(\cos A)) & Write the original equation. \\
(t^2=9^2+17^2-2(9)(17)\cos65°) & Substitute. \\
(t^2=81+289-306\cos65°) & Simplify. \\
(t^2\approx240.68) & Compute using technology. \\
(t\approx\sqrt{240.68}) & Take the square root of both sides. \\
(t\approx15.5) & Solve. \\
\end{tabular}
\te{COMMON ERROR Make sure to assign the correct values for a, b, and c based on their relationship to the given angle.}
\answer{To the nearest tenth, (t\approx15.5).}
\end{example}
\begin{tryit}{5}
a. In (\triangle JKL), (j=15), (k=13), and (l=12). What is (m\angle J)?
b. In (\triangle ABC), (a=11), (b=17), and (m\angle C=42°). What is c?
\answer{a. 73.6° b. about 11.5}
\end{tryit}
\begin{te-addendum}{5}{ElicitEvidence}
Elicit and Use Evidence of Student Thinking ETP
Q: When writing the Law of Cosines equation in part (a), how do you know which side to designate as a?
\answer{Because you are looking for (m\angle J), you need to substitute cos J for cos A, and therefore designate j as a.}
\end{te-addendum}
\begin{te-addendum}{5}{RtISupport}
\textbf{Support Struggling Students}
USE WITH EXAMPLE 5 Students may struggle with setting up the equation. Have them practice assigning the correct values for a, b, and c and setting up the equation for the Law of Cosines.
\item Use the Law of Cosines, (a^2=b^2+c^2-2bc(\cos A)) for (\triangle DEF).
% FIG 33 213 1154 483 1280
\placeholder{diagram}{Black triangle DEF, horizontal base D left F right 18 red; E above near right, DE 15, EF d; D angle 23° red.}
Q: Which angle corresponds to (\angle A) in the Law of Cosines?
\answer{(\angle D)}
Q: Which values correspond to a, b, and c in the Law of Cosines?
\answer{Side a corresponds to d, side b corresponds to 18 and side c corresponds to 15.}
Q: How do you set up the equation for the Law of Cosines?
\answer{Write the equation (a^2=b^2+c^2-2bc(\cos A)), then replace the variables with their corresponding values; (d^2=18^2+15^2-2(18)(15)(\cos23°)).}
\end{te-addendum}
% Source page 34: 409 Apply the Law of Cosines
\begin{te-addendum}{6}{GeneralizeETP}
\textbf{Apply the Law of Cosines}
Use and Connect Mathematical Representations ETP
Q: Why do you use the Law of Cosines and the Law of Sines to solve?
\answer{It is necessary because the flagpole is not perpendicular to the ground.}
Q: How do you know which angle to designate as (\angle A)?
\answer{Use the included angle that is given as (\angle A).}
Q: Why do you have to use the Law of Cosines first and then the Law of Sines to find the angle measure?
\answer{In order to use the Law of Sines, you need to know at least one side length and angle that are opposite each other. Since none are given in the problem, you first use the Law of Cosines. Once you have determined another value in the triangle, you have enough information to use the Law of Sines.}
\te{Printed teacher questions say Law of Sines second; the worked student example uses Law of Cosines twice. Either can find the remaining angle with suitable care. Examples are available at SavvasRealize.com.}
\end{te-addendum}
\begin{example}{6}
\textbf{Apply the Law of Cosines}
\te{APPLICATION}
A support wire is attached to the top of a flagpole that leans to the left. The wire is anchored to the ground 1.2 m to the right of the pole. The wire forms a 68° angle with the ground. What angle does the wire form with the top of the flagpole? Round to the nearest degree.
% FIG 34 960 244 1149 372
\placeholder{illustration}{US flag atop pole leaning left on green ground against blue sky. Wire from top to ground right of pole labeled 5.3 m. Horizontal dimension between pole base and anchor 1.2 m. Arc at wire anchor 68°.}
\te{STUDY TIP It is necessary to use the Law of Cosines and the Law of Sines because the flagpole is not perpendicular to the ground; if it were, you could use trigonometric ratios.}
Since the lengths of two sides and the measure of the included angle are given, use the Law of Cosines to find the length of the flagpole.
(a^2=b^2+c^2-2bc(\cos A))
(a^2=1.2^2+5.3^2-2(1.2)(5.3)\cos68°)
(a^2=1.44+28.09-12.72\cos68°)
(a^2\approx24.77)
(a\approx\pm\sqrt{24.77})
(a\approx4.98\text{ m})
Now that all side lengths are known, use the Law of Cosines to find the missing angle.
(b^2=a^2+c^2-2ac(\cos B))
(1.2^2=4.98^2+5.3^2-2(4.98)(5.3)(\cos B))
(1.44=24.8004+28.09-52.788(\cos B))
(\frac{-51.4504}{-52.788}=\cos B)
(\cos^{-1}(0.9747)\approx B)
(12.9°\approx B)
\te{USE APPROPRIATE TOOLS When finding angles, it is good practice to use the Law of Cosines whenever possible so you do not have to consider ambiguous cases.}
\answer{The angle formed by the wire and the flagpole is about 13°.}
\end{example}
\begin{tryit}{6}
A bike race follows a triangular path, represented by triangle ABC. If A is the starting point and the measure of the angle at point B is 70°, what is the measure of the angle formed by path BC and path CA?
% FIG 34 1083 707 1151 806
\placeholder{diagram}{Tall narrow black triangle ABC, A lower left, B lower right, C top. AB red 2 km, BC red 4 km, B angle red 70°.}
\answer{about 29.5°}
\end{tryit}
\begin{te-addendum}{6}{ElicitEvidence}
Elicit and Use Evidence of Student Thinking ETP
Q: What do you need to solve for first when using the Law of Cosines?
\answer{AC is needed in order to find (m\angle C).}
\end{te-addendum}
\begin{te-addendum}{6}{HabitsOfMind}
Use with EXAMPLES 4, 5, \& 6
Use Appropriate Tools Given the lengths of two sides of a triangle, and the measure of one angle, how do you know whether to use the Law of Sines or the Law of Cosines?
\answer{If the angle whose measure is given is the included angle between the two given sides, then use the Law of Cosines. If the angle is opposite one of the known sides, then use the Law of Sines.}
\end{te-addendum}
% Source page 35: 410 Concept Summary
\begin{te-addendum}{ConceptSummary}{PurposefulQuestions}
\textbf{CONCEPT SUMMARY Law of Sines and Law of Cosines}
Q: Why are the Law of Sines and the Law of Cosines helpful when finding missing side lengths and angle measures?
\answer{These laws allow you to find missing side lengths and angle measurements in acute and obtuse triangles.}
\te{Presentation. Practice. Concept Summary, Do You Understand, and Do You Know How are available at SavvasRealize.com.}
\end{te-addendum}
\begin{te-addendum}{ConceptSummary}{CommonError}
Exercise 9 Some students may try to use the Law of Sines. Ask students to identify the angles and sides that are opposite each other. Then ask whether they are given both values for any of the pairs. Show them that in order to use the Law of Sines they must know both values for at least one of these pairs.
\end{te-addendum}
\begin{concept-summary}
\textbf{CONCEPT SUMMARY Law of Sines and Law of Cosines}
\begin{tabular}{lll}
 & Law of Sines & Law of Cosines \\
ALGEBRA & (\frac{\sin A}{a}=\frac{\sin B}{b}=\frac{\sin C}{c}) & (a^2=b^2+c^2-2bc(\cos A)) \\
 & & (b^2=a^2+c^2-2ac(\cos B)) \\
 & & (c^2=a^2+b^2-2ab(\cos C)) \\
\end{tabular}
% FIG 35 744 267 866 321
\placeholder{diagram}{Black triangle ABC, horizontal base A left C right labeled red b, B upper; AB red c and BC red a.}
% FIG 35 944 267 1070 321
\placeholder{diagram}{Black triangle ABC, horizontal base A left B right red c, C upper; AC red b and CB red a.}
NUMBERS. Law of Sines
% FIG 35 744 387 862 441
\placeholder{diagram}{Black triangle ACB, base A left B right red 49, C above, angle C red 115°, AC b, CB 16.}
Find the measure of (\angle A).
(\frac{\sin A}{a}=\frac{\sin C}{c}\text{ so }\frac{\sin A}{16}=\frac{\sin115°}{49})
(\frac{\sin A}{16}=\frac{\sin115°}{49})
(\sin A\approx0.2959)
(A\approx\sin^{-1}0.2959;m\angle A\approx17.2°)
NUMBERS. Law of Cosines
% FIG 35 945 386 1070 455
\placeholder{diagram}{Black triangle ACB, base A left B right red sqrt2, C above, AC b, CB 5, angle B 45° red.}
Find b.
(b^2=a^2+c^2-2ac(\cos B))
(b^2=5^2+(\sqrt2)^2-2(5)(\sqrt2)(\cos45°))
(b^2=25+2-10(\sqrt2)(\frac{\sqrt2}{2}))
(b^2=17;b\approx4.1)
\textbf{Do You UNDERSTAND?}
1. ESSENTIAL QUESTION How can you use the sine and cosine functions with non-right triangles?
\answer{You can use the Law of Sines and the Law of Cosines in non-right triangles to find the measures of sides and angles when you have a side and its corresponding angle as well as another side or angle.}
2. Error Analysis Alejandro said the Law of Sines always gives one answer. Explain and correct Alejandro's error.
\answer{The Law of Sines can give zero, one, or two answers. When solving for an unknown angle, if the solution is an acute angle, then you will need to use that answer to find a second angle. If there is no solution, the Law of Sines gives no valid solutions.}
\te{Printed wording implies every acute solution yields a second triangle; likely the supplementary candidate must also satisfy the triangle angle sum.}
3. Construct Arguments Consider the Law of Cosines as (a^2=b^2+c^2-2bc(\cos A)). Explain why the negative square root of a is not a valid solution.
\answer{The negative square root is not valid because a is a side length and side lengths cannot be negative.}
\te{Printed ``negative square root of a''; likely negative square root when solving for a.}
4. Reason In what situations do you use the Law of Sines? Law of Cosines?
\answer{The Law of Sines is used to solve for a side or angle measure when you have a side and angle pair (a and A, for example), and one more side or angle. The Law of Cosines is used to find an angle when you have all three sides or to find a side when you have the other two sides and the included angle.}
\textbf{Do You KNOW HOW?}
Use the Law of Sines or the Law of Cosines to find the indicated measure in (\triangle ABC).
5. (m\angle A=50°), (a=4.5), (b=3.8); find (m\angle B).
\answer{about 40.3°}
6. (m\angle A=72°), (a=61), (c=58); find (m\angle C).
\answer{about 64.7°}
7. (m\angle A=18°), (m\angle C=75°), (c=101); find a.
\answer{about 32.3}
8. (m\angle B=112°), (m\angle C=20°), (c=1.6); find b.
\answer{about 4.3}
9. (m\angle C=45°), (a=15), (b=8); find c.
\answer{about 10.9}
10. (m\angle A=82°), (b=2.5), (c=6.8); find a.
\answer{about 6.9}
11. (a=14), (b=12), (c=5.8); find (m\angle A).
\answer{about 97.6°}
\end{concept-summary}
% Source page 36: 411 Practice & Problem Solving
\begin{te-addendum}{Practice}{GeneralizeETP}
\textbf{PRACTICE \& PROBLEM SOLVING}
Practice \& Problem Solving and video tutorials are available at SavvasRealize.com.
Lesson Practice You may opt to have students complete the automatically scored Practice and Problem Solving items online powered by MathXL for School.
Choose from: Lesson Practice; Adaptive Practice.
You may also take advantage of the bank of exercises for assigning additional practice.
\textbf{Assignment Guide}
\begin{tabular}{ll}
On-Level & Advanced \\
12--26,28--37 & 12--17,19--37 \\
\end{tabular}
\textbf{Engage Through Student Choice}
Promote student agency by allowing students to choose practice items. You may structure this choice in many ways.
For example:
Assign each section a point value. Students choose at least one item from each section and items chosen should have a minimum of 20 total points.
\begin{tabular}{ll}
Understand, Apply & 2 points each \\
Practice & 1 point each \\
Assessment Practice & 1 point each \\
Performance Task & 3 points \\
\end{tabular}
\te{Scan the QR code to access a video tutorial.}
\end{te-addendum}
\begin{item-analysis}
\begin{tabular}{lll}
\toprule
Example & Items & DOK \\
\midrule
1 & 18,19 & 2 \\
2 & 20,21,36 & 1 \\
2 & 12,32 & 2 \\
3 & 13 & 1 \\
3 & 17,22--24,35 & 2 \\
4 & 16 & 2 \\
4 & 25 & 3 \\
5 & 26,27 & 1 \\
5 & 14,15,30,31 & 2 \\
6 & 28,29 & 2 \\
6 & 33,34,37 & 3 \\
\bottomrule
\end{tabular}
\end{item-analysis}
\begin{practice}{12}{2}
UNDERSTAND. Construct Arguments Lourdes said that you can use the Law of Sines if you have any two angles and any side, or any two sides and any angle. Is Lourdes correct? Explain your reasoning.
\answer{No, Lourdes is not correct. To use the Law of Sines, you need a corresponding side and angle, plus one other side or angle. If you have sides a and b, and angle C, for example, it would not work.}
\end{practice}
\begin{practice}{13}{1}
Generalize What combination of sides and/or angles in a triangle leads to the ambiguous case?
\answer{Angle-Side-Side}
\end{practice}
\begin{practice}{14}{2}
Error Analysis Describe and correct the error a student made in using the Law of Cosines to find a.
(a^2=14^2+9^2-2(14)(9)(\cos140))
(a^2=196+81-252(-0.766))
(a^2=277-193.03)
(a^2=83.97)
(a\approx9.16)
% FIG 36 572 476 791 602
\placeholder{illustration}{Pale gray student work note with curled lower right corner and red X, containing the five incorrect calculation lines transcribed above.}
\answer{Since cos 140° is negative, the student should have added 193.03 to 277 to get (a^2=470.03). Then (a=21.68).}
\end{practice}
\begin{practice}{15}{2}
Communicate Precisely Two students are solving for d in the triangle shown, using the Law of Cosines. One says the answer is 20 in., and the other says the answer is 20.3 in. Is either student incorrect? Explain your reasoning.
% FIG 36 567 701 705 790
\placeholder{diagram}{Black triangle with left angle red 46°, upper vertex, lower right vertex; side from left to upper 17 in., left to lower right 28 in., upper to lower right d.}
\answer{Neither student is incorrect; they reported their answers to different degrees of precision. Since the given values are in integers, 20 in. is probably precise enough.}
\end{practice}
\begin{practice}{16}{2}
Look for Relationships Show that the Law of Cosines is equivalent to the Pythagorean Theorem when the given angle is 90°.
% FIG 36 565 850 694 930
\placeholder{diagram}{Black right triangle with horizontal bottom leg red b, vertical left leg red a, descending hypotenuse red c; lower left C with red right-angle square.}
\answer{(c^2=a^2+b^2-2ab(\cos90)); (c^2=a^2+b^2-2ab(0)); (c^2=a^2+b^2-0); (c^2=a^2+b^2). This is the Pythagorean Theorem.}
\end{practice}
\begin{practice}{17}{2}
Higher Order Thinking The ambiguous case only causes a problem when the given angle is acute, not when an obtuse angle is given. Explain why.
\answer{Inverse sine gives an acute angle, and the supplement to that angle will always be obtuse. This causes a problem when the known angle is acute because both the acute angle and its obtuse supplement could be valid options for the missing angle. However, if the known angle is obtuse, only the acute solution is valid because a triangle cannot contain two obtuse angles.}
\end{practice}
\begin{practice}{18}{2}
PRACTICE. How can you derive the Law of Sines for the given angles? SEE EXAMPLE 1
E and F
% FIG 36 939 312 1064 378
\placeholder{diagram}{Black obtuse triangle EFG, E upper left, F below and slightly right, G far right at same height as F; horizontal FG.}
\answer{Draw the altitude from G, label its length x. (\sin E=\frac xf) and (\sin F=\frac xe); (f\sin E=x) and (e\sin F=x); (f\sin E=e\sin F); (\frac{\sin E}{e}=\frac{\sin F}{f}).}
\end{practice}
\begin{practice}{19}{2}
How can you derive the Law of Sines for the given angles? SEE EXAMPLE 1
F and G
\te{Uses the triangle shown with Exercise 18.}
\answer{Draw the altitude from E, label its length x. (\sin F=\frac xg) and (\sin G=\frac xf); (g\sin F=x) and (f\sin G=x); (g\sin F=f\sin G); (\frac{\sin F}{f}=\frac{\sin G}{g}).}
\te{Printed proof keys for 18 and 19 omit the supplementary-angle step needed for obtuse F; likely use sin(180°-F)=sin F.}
\end{practice}
\begin{practice}{20}{1}
Use the Law of Sines to solve. SEE EXAMPLE 2
In (\triangle HJK), (m\angle J=122°), (j=17), and (k=8). What is (m\angle K)?
\answer{about 23.5°}
\end{practice}
\begin{practice}{21}{1}
Use the Law of Sines to solve. SEE EXAMPLE 2
In (\triangle RST), (m\angle R=45°), (m\angle S=19°), and (r=15). What is s?
\answer{about 6.9}
\end{practice}
\begin{practice}{22}{2}
Find the number of possible triangles for each set of measures. Then find the angle measure(s). SEE EXAMPLE 3
In (\triangle WXY), (m\angle X=104°), (x=7), and (y=2). Find (m\angle Y).
\answer{one: (m\angle Y\approx16.1°)}
\end{practice}
\begin{practice}{23}{2}
Find the number of possible triangles for each set of measures. Then find the angle measure(s). SEE EXAMPLE 3
In (\triangle DEF), (m\angle D=28°), (d=8), and (e=15). Find (m\angle E).
\answer{two: (m\angle E\approx61.7°) or 118.3°}
\end{practice}
\begin{practice}{24}{2}
Find the number of possible triangles for each set of measures. Then find the angle measure(s). SEE EXAMPLE 3
In (\triangle RST), (m\angle R=30°), (r=14), and (t=32). Find (m\angle T).
\answer{zero}
\end{practice}
\begin{practice}{25}{3}
How can you derive the Law of Cosines for obtuse angle K? SEE EXAMPLE 4
% FIG 36 877 694 959 756
\placeholder{diagram}{Black triangle JKL, horizontal top JK, J left K right, L below and right of K; obtuse angle at K.}
\answer{Sample: Student draws altitude outside triangle from vertex L. (\cos(180°-K)=-\cos K=\frac xj); (x=-j\cos K); (x^2+h^2=j^2) and ((l+x)^2+h^2=k^2); (l^2+2lx+x^2+h^2=k^2); (k^2=l^2+2l(-j\cos K)+j^2); (k^2=l^2+j^2-2lj(\cos K)).}
\end{practice}
\begin{practice}{26}{1}
In (\triangle PQR), find (m\angle P). SEE EXAMPLE 5
(p=5), (q=8), (r=9)
\answer{about 33.6°}
\end{practice}
\begin{practice}{27}{1}
In (\triangle PQR), find (m\angle P). SEE EXAMPLE 5
(p=14), (q=6), (r=12)
\answer{about 96.4°}
\end{practice}
\begin{practice}{28}{2}
What is the measure of angle Z? SEE EXAMPLE 6
% FIG 36 868 852 1017 935
\placeholder{diagram}{Black triangle XYZ, horizontal base Y left Z right, X above near right. Red XY 54, YZ 42, angle Y 25°.}
\answer{about 106.9°}
\end{practice}
\begin{practice}{29}{2}
What is the measure of angle Z? SEE EXAMPLE 6
% FIG 36 1042 855 1139 980
\placeholder{diagram}{Black triangle XYZ, X left, Y upper right, Z bottom. Red XY 7.7, YZ 9, angle Y 34°.}
\answer{about 58.7°}
\end{practice}
% Source page 37: 412 Apply and Assessment Practice
\begin{practice}{30}{2}
APPLY. Model With Mathematics The head sail for Mekelle's sailboat is a triangle with the three sides having lengths of 24 ft, 23 ft, and 12 ft. What is the measure of the sail's greatest angle?
\answer{about 79.9°}
\end{practice}
\begin{practice}{31}{2}
Use Structure An art sculpture is made of rotated scalene triangles, as shown. The triangles are all congruent. What is the length of the longest side of each triangle?
% FIG 37 626 368 778 527
\placeholder{illustration}{Gray sculpture of three intersecting rotated open triangular frames on a base. Foreground triangle has bottom dimension 40 in., right side dimension 78 in., included lower right angle 100°.}
\answer{about 93.6 in.}
\end{practice}
\begin{practice}{32}{2}
Make Sense and Persevere The course for a race follows three roads as shown. How far do the runners travel along Jappa Road?
% FIG 37 517 587 776 730
\placeholder{map}{Three gray roads enclose a black triangular route: Jappa Rd. horizontal top, Ash St. descends right from upper left to bottom, Carson St. ascends right from bottom to upper right. Ash side 38 yd, bottom angle 80°, upper right angle 60°. Trees, buildings and parked cars surround roads.}
\answer{about 43.2 yd}
\end{practice}
\begin{practice}{33}{3}
Model With Mathematics Noemi throws a ball to Roberto, who is 6 m away. When Roberto catches the ball, he turns 50°, and then throws the ball 7 m to Shandra. What angle does Shandra turn to throw back to Noemi?
\answer{about 55.6°}
\end{practice}
\begin{practice}{34}{3}
Make Sense and Persevere Tamika parked her car and walked 300 yd down a path. She then made a 135° turn onto a new path. She walked another 40 yd along a river to her fishing spot. If Tamika turns to face the direction of her car, what angle does she need to turn?
\answer{about 40.1°}
\end{practice}
\begin{practice}{35}{2}
ASSESSMENT PRACTICE. In (\triangle EFG), (m\angle E=35°), (e=5.8), and (f=10). Choose Yes or No to tell whether each is a possible value for (m\angle F).
\begin{tabular}{lll}
 & Yes & No \\
There are no possible value & blank & blank \\
6.2 & blank & blank \\
60.3 & blank & blank \\
81.5 & blank & blank \\
98.5 & blank & blank \\
119.7 & blank & blank \\
\end{tabular}
\answer{No; No; No; Yes; Yes; No.}
\te{Printed ``There are no possible value''; likely ``values''.}
\end{practice}
\begin{practice}{36}{1}
SAT/ACT In (\triangle ABC), (a=29.7), (b=48.5), and (B=92°). What is (m\angle A)?
A. There is no possible value.
B. 37.7°
C. 56.3°
D. 123.7°
E. 142.3°
\answer{B}
\end{practice}
\begin{practice}{37}{3}
Performance Task Teo is standing 80 yd from the base of a Mayan pyramid. The side of the pyramid is 100 ft from base to peak.
% FIG 37 814 700 1085 858
\placeholder{photo}{Mayan pyramid under cloudy blue sky. Black triangle overlays Teo lower left, pyramid base lower right, peak upper right. Horizontal base 80 yd, pyramid side 100 ft, sight line x; Teo angle 20°, peak angle P.}
Part A How many feet from the base of the pyramid is Teo?
\answer{Part A 240 ft}
Part B What is the measure of angle P formed by the side of the pyramid and Teo's line of sight?
\answer{Part B about 55.2°}
Part C What is the distance in a straight line from Teo to the peak x?
\answer{Part C about 282.6 ft}
\end{practice}
% Source page 38: 412A Assess & Differentiate
\begin{te-addendum}{Assess}{RtISupport}
\textbf{LESSON QUIZ}
Use the Lesson Quiz to assess students' understanding of the mathematics in the lesson.
Students can take the Lesson Quiz online or you can download a printable copy from SavvasRealize.com. The Lesson Quiz is also available in the Assessment Sourcebook.
\textbf{Item Analysis}
\begin{tabular}{lll}
Item & DOK & Skills Review and Practice \\
1 & 2 & G96 \\
2 & 1 & G96 \\
3 & 2 & G97 \\
4 & 1 & G97 \\
5 & 2 & G96 \\
\end{tabular}
Use the student scores on the Lesson Quiz to prescribe differentiated assignments.
If students take the Lesson Quiz online, it will be automatically scored and appropriate differentiated practice will be assigned based on student performance.
\begin{tabular}{lll}
Intervention & 0--3 points & Reteach to Build Understanding; Mathematical Literacy and Vocabulary; Lesson Virtual Nerd videos \\
On-Level & 4 points & Enrichment \\
Advanced & 5 points & Enrichment \\
\end{tabular}
\te{Assessment. Lesson Quiz is available at SavvasRealize.com.}
\end{te-addendum}
\begin{te-addendum}{Assess}{GeneralizeETP}
\textbf{8-2 Lesson Quiz: Law of Sines and Law of Cosines}
\te{ASSESSMENT RESOURCES. Name blank. enVision Algebra 2. SavvasRealize.com. Copyright© Savvas Learning Company LLC. All Rights Reserved. enVision Algebra 2 Assessment Sourcebook.}
1. In (\triangle PQR), the length of the altitude from Q is x. Complete the derivation of the Law of Sines for (\triangle PQR).
Choices: (\frac xr); (r\sin P); (\frac xp); (p\sin R); (r\sin P=p\sin R); (\frac{\sin P}{p}=\frac{\sin R}{r}).
% FIG 38 721 408 849 465
\placeholder{diagram}{Black triangle PQR, horizontal base P left R right q; Q above near right, PQ r, QR p. Vertical altitude x from Q to base with dashed right-angle square at foot.}
(\sin P=\text{blank}) so (x=\text{blank})
(\sin R=\text{blank}) so (x=\text{blank})
Setting the two expressions equal to each other gives the equation blank.
Divide both sides of the equation by pr to give the result: blank.
\answer{1. (\frac xr); (r\sin P); (\frac xp); (p\sin R); (r\sin P=p\sin R); (\frac{\sin P}{p}=\frac{\sin R}{r}).}
2. In (\triangle LMN), (m\angle L=108°), (m\angle M=25°), and (\ell=12). What is m?
A. about 0.19
B. about 5.3
C. about 20.8
D. 27
\answer{2. B}
3. Tia uses the Law of Cosines to find a missing length in (\triangle ABC). After Tia substitutes the values she knows into the equation and simplifies, she gets (-23.26\cos A=0). What does this tell you about (\triangle ABC)?
A. There are two possible triangles that match the information Tia is using.
B. The triangle is a right triangle.
C. There is no triangle that matches the information Tia is using.
D. The triangle is an obtuse triangle.
\answer{3. B}
4. What is the length y? Round to the nearest tenth.
% FIG 38 994 724 1110 786
\placeholder{diagram}{Black triangle XYZ with nearly horizontal upper side XY 10, X left Y right, Z below; ZY 7, XZ y, angle Y 32°.}
(y\approx\text{blank})
\answer{4. 5.5}
5. A path in a park forms a triangle as shown. What is the measure of angle A? Round to the nearest tenth of a degree.
% FIG 38 995 784 1112 838
\placeholder{diagram}{Black triangle ACB, horizontal base A left B right, C above left of center. AC 72 yd, CB 135 yd, angle C 119°.}
(m\angle A\approx\text{blank}°)
\answer{5. 40.7}
\end{te-addendum}
% Source page 39: 412B Differentiated Resources
\begin{te-addendum}{Assess}{GeneralizeETP}
\textbf{DIFFERENTIATED RESOURCES}
I=Intervention; O=On-Level; A=Advanced.
Reteach to Build Understanding I. Provides scaffolded reteaching for the lesson concepts. MathXL for School.
Additional Practice I O. Provides extra practice for each lesson. MathXL for School.
Enrichment O A. Presents engaging problems and activities that extend the lesson concepts. MathXL for School.
Mathematical Literacy and Vocabulary I O. Helps students develop and reinforce understanding of key terms and concepts.
\te{Resource sheets each print Name blank, enVision Algebra 2, SavvasRealize.com, lesson title Law of Sines and Law of Cosines, Copyright© Savvas Learning Company LLC. All Rights Reserved., and enVision Algebra 2 Teaching Resources.}
\end{te-addendum}
\begin{te-addendum}{Assess}{RtISupport}
\textbf{8-2 Reteach to Build Understanding}
1. Examine the Law of Sines by comparing the ratio of the sine of each angle to its corresponding side. Round your answer to the nearest thousandth.
% FIG 39 160 432 238 495
\placeholder{diagram}{Black triangle ABC, horizontal base A left B right dimension c=21, C above left. AC b=13, CB a=20, altitude h=12 to base with right-angle square, splitting base into 5 left and 16 right. Angles A 67.3°, B 36.9°, C 75.8° with arc and leader.}
\begin{tabular}{ll}
The Law of Sines & (\frac{\sin A}{a}=\frac{\sin B}{b}=\frac{\sin C}{c}) \\
(\frac{\sin A}{a}) & (\frac{\sin\uncertain{67.3°}{67.4°}}{20}=\answer{0.046}) \\
(\frac{\sin B}{b}) & (\frac{\sin36.9°}{13}=\answer{0.046}) \\
(\frac{\sin C}{c}) & (\frac{\sin75.8°}{21}=\answer{0.046}) \\
\end{tabular}
2. Emery finds the measure of (\angle A) in the triangle. What mistake does she make? What is the correct measure?
% FIG 39 160 534 294 585
\placeholder{diagram}{Black obtuse triangle CBA, C upper left, B below and slightly right, A far right, BA horizontal 19, BC 6, CA 23, angle B 122°.}
(\frac{\sin A}{a}=\frac{\sin B}{b})
(\frac{\sin A}{19}=\frac{\sin122°}{23})
(A=\sin^{-1}(\frac{19\sin122°}{23})\approx43.9°)
\answer{Emery substituted the wrong side length. The correct answer is 12.6°.}
3. Use the Law of Cosines to find the value of a.
% FIG 39 317 608 408 673
\placeholder{diagram}{Black triangle ABC, horizontal base AC 10, A lower left C lower right B above; AB 8, BC a, angle A 60°.}
(a^2=b^2+c^2-2bc(\cos A))
(a^2=10^2+8^2-2(\answer{10})(\answer{8})(\cos\answer{60}°))
(a^2=\answer{10}^2+8^2-2(10)(8)(\answer{0.5}))
(a^2=\answer{100}+\answer{64}-2(10)(8)(0.5))
(a^2=100+64-\answer{80})
(a^2=\answer{164}-80=\answer{84})
(a\approx9.165)
\end{te-addendum}
\begin{te-addendum}{Assess}{GeneralizeETP}
\textbf{8-2 Additional Practice}
1. How can you derive the Law of Sines for angles A and C?
% FIG 39 531 428 622 463
\placeholder{diagram}{Black triangle ABC, horizontal base A left C right, B above nearer A, vertical altitude x from B to AC.}
\answer{(\sin A=\frac xc) or (x=c\sin A) and (\sin C=\frac xa) or (x=a\sin C); (c\sin A=a\sin C\Rightarrow\frac{\sin A}{a}=\frac{\sin C}{c}).}
2. In (\triangle MNO), (m\angle M=135°), (m=18), and (n=14). Find (m\angle O). Round your answer to the nearest tenth.
\answer{11.6°}
3. In (\triangle ABC), (m\angle A=60°), (a=8), and (b=6). Find (m\angle B). Round your answer to the nearest tenth.
\answer{40.5°}
4. Describe and correct the error a student made in using the Law of Cosines to solve for b in (\triangle ABC) where (m\angle B=120°), (a=16), and (c=14).
\answer{The student did not multiply ac cos B by 2.}
\begin{tabular}{ll}
Student work & Corrected work \\
(b^2=16^2+14^2-(16)(14)(\cos120°)) & (b^2=16^2+14^2-\answer{2}(16)(14)(\cos120°)) \\
(b^2=256+196-(224)(-0.5)) & (b^2=256+196-(\answer{448})(-0.5)) \\
(b^2=256+196+112=564) & (b^2=256+196+\answer{224}=676) \\
(b\approx23.7) & (b=\answer{26}) \\
\end{tabular}
5. The triangle illustrates the side view of a roof truss with edge-lengths of 11 ft and 23 ft. The angle between the edges is 55°. What is the length of AC to the nearest foot?
% FIG 39 529 672 619 711
\placeholder{diagram}{Black triangle ABC, A lower left C lower right B upper left; horizontal AC, near-vertical AB 11 ft, BC 23 ft, angle B 55°.}
\answer{19 ft}
6. Dyani throws a ball to Edgar, who is 8 m away. When Edgar catches the ball, he turns 55°, and then throws the ball 9 m to Hana. What angle does Hana turn to throw the ball to Dyani? Round to the nearest tenth.
\answer{56.1°}
\end{te-addendum}
\begin{te-addendum}{Assess}{RtIExtend}
\textbf{8-2 Enrichment}
Triangles ABC and MBC have common side BC. The side-lengths of the triangle ABC are denoted by a, b, and c, and the non-common sides of the triangle MBC are denoted by m, and n.
The Law of Cosines takes the following forms:
(a^2=2b^2-2bc\cos(60))
(a^2=2c^2-2bc\cos(60))
(a^2=2m^2-2mn\cos(20))
(a^2=2n^2-2mn\cos(20))
1. Which statements are true? Draw a line through the statements that are not true.
\item Triangle ABC is isosceles.
\item Triangle MBC is isosceles.
\item Triangle ABC equilateral.
\item Triangle MBC is equilateral.
\item Triangle ABC is scalene.
\item Triangle MBC is scalene.
\item The difference between angles M and A is 40 degrees.
\item The non-common sides of MBC are longer than the non-common sides of ABC.
\item The non-common sides of MBC cannot be compared with sides of ABC.
\answer{Cross out statements 4,5,6,7,9. Leave statements 1,2,3,8.}
\te{Printed statement 7 is crossed out, but the given cosine arguments imply angles A=60° and M=20°; likely the 40-degree difference is true.}
2. Right triangle ABC shares side AC with the scalene triangle ACD. Prove that
(\cos D=\frac{p^2+q^2-a^2-b^2}{2pq}).
% FIG 39 855 632 956 700
\placeholder{diagram}{Black quadrilateral ABCD with A upper left B lower left C lower right D upper right. AB vertical c, BC horizontal a, right-angle square at B; diagonal AC b; AD p, DC q.}
\answer{(b^2=p^2+q^2-2pq\cos D) and (b^2=a^2+c^2); (a^2+c^2=p^2+q^2-2pq\cos D); (a^2+c^2-p^2-q^2=-2pq\cos D); (\frac{a^2+c^2-p^2-q^2}{-2pq}=\cos D); (\frac{p^2+q^2-a^2-c^2}{2pq}=\cos D).}
\te{Printed requested formula subtracts b squared; the proof subtracts c squared, which is the likely intended formula.}
\end{te-addendum}
\begin{te-addendum}{Assess}{ELLAddendum}
\textbf{8-2 Mathematical Literacy and Vocabulary}
Side-lengths and angle-measures of (\triangle ABC) are given. Some of the measures are expressed in terms of a, b, and c.
% FIG 39 341 995 421 1037
\placeholder{diagram}{Black triangle ABC, horizontal base B left C right labeled 49, A above near left. AB a-c, AC a-b; angles A 100°, B 60°, C 20°.}
Teddy wrote the steps to find the values of (a-b), (a-c), and (c-b) on note cards, but they got mixed up. Use the cards to list the steps in the correct order. Start with the side AB.
Cards:
\item (\frac{0.98}{49}=\frac{0.34}{a-c})
\item ((a-b)-(a-c)=43.5-17)
\item (\frac{\sin100}{49}=\frac{\sin20}{a-c})
\item (a-c=17)
\item (\frac{0.98}{49}=\frac{0.34}{a-c})
\item (a-b=43.5)
\item (c-b=26.5)
\item (\frac{\sin100}{49}=\frac{\sin60}{a-b})
\answer{1. (\frac{\sin100}{49}=\frac{\sin20}{a-c})}
\answer{2. (\frac{0.98}{49}=\frac{0.34}{a-c})}
\answer{3. (a-c=17)}
\answer{4. (\frac{\sin100}{49}=\frac{\sin60}{a-b})}
\answer{5. (\frac{0.98}{49}=\frac{0.87}{a-b})}
\answer{6. (a-b=43.5)}
\answer{7. ((a-b)-(a-c)=43.5-17)}
\answer{8. (c-b=26.5)}
\te{Printed cards repeat the 0.34/(a-c) equation; likely one should be the 0.87/(a-b) equation used in answer 5.}
\end{te-addendum}
\begin{te-addendum}{Assess}{GeneralizeETP}
\textbf{Digital Differentiated Resources}
The Reteach to Build Understanding, Additional Practice, and Enrichment activities are available as digital assignments. These activities are automatically assigned when students complete the lesson quiz online and are automatically scored.
Solve the system of equations by graphing.
(y=3x+4)
(y=-2x+4)
Use the graphing tool to graph the system.
Click to enlarge graph.
Click the graph, choose a tool in the palette and follow the instructions to create your graph.
% FIG 39 584 978 1035 1300
\placeholder{illustration}{Black landscape tablet with home button left, white MathXL screen. System equations left, empty coordinate grid upper right, arrowed x,y axes, unit grid and labels -10,-8,...,10. Question Help menu Help Me Solve This, View an Example, Video, Textbook, Glossary, Math Tools, Print. Footer Review progress, Question 5 of 14, Go, Back, Next, Check Answer, Clear All; 1 part remaining.}
\end{te-addendum}

\end{document}
